% =====================================================================
%  Lisans Bitirme Tezi --- Sunum (Beamer)
%  Yüz Tanıma Tabanlı Sürücü Güvenlik ve Acil Durum Bildirim Sistemi
%  Hedef süre: ~15 dakika defans + sergi sunumu
%
%  Derleme:  latexmk -pdf main.tex
% =====================================================================
\documentclass[aspectratio=169,12pt]{beamer}

\usepackage[utf8]{inputenc}
\usepackage[T1]{fontenc}
\usepackage[turkish]{babel}
\usepackage{lmodern}
\usepackage{microtype}
\usepackage{booktabs}
\usepackage{tabularx}
\usepackage{tikz}
\usetikzlibrary{shapes.geometric, arrows.meta, positioning, fit, backgrounds}

% --- Renk paleti (poster + tez ile uyumlu) ------------------------------
\definecolor{ink}{HTML}{11162A}
\definecolor{accent}{HTML}{8C631A}
\definecolor{paper}{HTML}{F1ECE0}
\definecolor{mute}{HTML}{5B6075}
\definecolor{signal}{HTML}{9C2A21}

% --- Beamer teması: minimal, akademik -----------------------------------
\usetheme{default}
\usecolortheme{seahorse}
\useinnertheme{circles}
\useoutertheme{infolines}

\setbeamercolor{normal text}{fg=ink, bg=paper}
\setbeamercolor{title}{fg=ink}
\setbeamercolor{frametitle}{fg=ink, bg=paper}
\setbeamercolor{structure}{fg=accent}
\setbeamercolor{block title}{fg=paper, bg=ink}
\setbeamercolor{block body}{bg=ink!4, fg=ink}
\setbeamercolor{item}{fg=accent}
\setbeamercolor{alerted text}{fg=signal}

\setbeamerfont{frametitle}{series=\bfseries, size=\large}
\setbeamerfont{title}{series=\bfseries, size=\LARGE}

\setbeamertemplate{navigation symbols}{}    % alt köşedeki navigasyonu sil
\setbeamertemplate{footline}[frame number]

% --- Başlık metadata -----------------------------------------------------
\title[Sürücü Güvenliği]{Yüz Tanıma Tabanlı Sürücü Güvenlik ve\\Acil Durum Bildirim Sistemi}
\subtitle{Driver Safety and Emergency Notification System Based on Facial Recognition}
\author[E.\ Ağaoğlu]{Ekin Ağaoğlu \\ \small Danışman: Prof.\ Dr.\ Aydoğan Savran}
\institute[Ege Ü.]{Ege Üniversitesi \\ Elektrik-Elektronik Mühendisliği Bölümü}
\date{Lisans Bitirme Projesi · 2025--2026}

% =========================================================================
\begin{document}

% --------------------------------------------- 1. Kapak --------------------
\begin{frame}[plain]
    \titlepage
\end{frame}

% --------------------------------------------- 2. Gündem ------------------
\begin{frame}{Sunum Akışı}
    \begin{enumerate}[<+->]
        \item Problem ve motivasyon
        \item Literatür ve tasarım kararları
        \item Sistem mimarisi
        \item Donanım ve yazılım gerçeklemesi
        \item Sürücü uygulaması (iPhone) + kullanıcı yönetimi
        \item Ölçüm planı ve sonuçlar
        \item Sonuç ve gelecek çalışmalar
    \end{enumerate}
\end{frame}

% =========================================================================
% Bölüm 1 — PROBLEM
% =========================================================================
\section{Problem ve Motivasyon}

\begin{frame}{Problem}
    Ticari taksilerde sürücü güvenliği:
    \begin{itemize}
        \item Yolcu kimliği sürüş öncesinde \alert{doğrulanmaz}.
        \item Mevcut çözümler (panik buton, GPS) çoğunlukla
              \alert{olay sonrası} bilgi sağlar.
        \item Sürücünün aktif eylem gerektiren sistemler dikkat dağıtıcıdır.
    \end{itemize}

    \vspace{0.4cm}
    İki temel ihtiyaç:
    \begin{enumerate}
        \item Caydırıcı bir \textbf{kimlik doğrulama} altyapısı.
        \item Sürücünün dikkatini bozmayan \textbf{otomatik bildirim}.
    \end{enumerate}
\end{frame}

\begin{frame}{Tasarım Hedefleri}
    \begin{block}{Başarı Ölçütleri}
        \begin{itemize}
            \item Uçtan uca tanıma gecikmesi: $\leq$ 7~s (10-frame burst)
            \item Eşit Hata Oranı (EER): $\leq$ \%5
            \item Aydınlatma: 100--600~lx · Mesafe: 30--120~cm
            \item Acil çağrı dialer'ının açılma süresi: $\leq$ 3~s
            \item Ek donanım maliyeti: $\leq$ 500~TL
        \end{itemize}
    \end{block}
\end{frame}

% =========================================================================
% Bölüm 2 — LİTERATÜR ve TASARIM KARARLARI
% =========================================================================
\section{Literatür ve Kararlar}

\begin{frame}{Yüz Tanıma --- Yöntem Seçimi}
    \begin{itemize}
        \item \textbf{ArcFace}: kosinüs benzerliği uzayında additive angular
              margin --- yüksek doğruluk \alert{[Deng 2019]}
        \item \textbf{InsightFace buffalo\_l}: açık kaynak, RetinaFace +
              ArcFace R100
        \item Kapalı kutu API'lar (AWS Rekognition vb.) reddedildi:
              KVKK belirsizliği, kontrolsüz model değişiklik riski
    \end{itemize}

    \vspace{0.3cm}
    \begin{exampleblock}{Kazanç}
        Model + eşik bizim kontrolümüzde $\rightarrow$ FAR/FRR ölçümü
        tezde dolgun savunma; KVKK uyumu doğal.
    \end{exampleblock}
\end{frame}

\begin{frame}{Tek Snapshot vs.\ Multi-Frame Burst}
    \begin{columns}
    \column{0.5\textwidth}
        \begin{block}{Tek-shot}
            \begin{itemize}
                \item Hareket bulanıklığına duyarlı
                \item Kırpılan yüze duyarlı
                \item Liveness sinyali yok
            \end{itemize}
        \end{block}
    \column{0.5\textwidth}
        \begin{block}{10-frame burst (seçildi)}
            \begin{itemize}
                \item 5 FPS × 2 s = 10 kare
                \item Centroid agregasyon $\rightarrow$ robust
                \item Cross-frame $\sigma$ $\rightarrow$ pasif canlılık
            \end{itemize}
        \end{block}
    \end{columns}

    \vspace{0.4cm}
    Maliyeti yalnızca $\sim$5--7~s gecikme --- yolcu binişi başına bir
    işlem, kullanıcı deneyimi için uygun.
\end{frame}

\begin{frame}{Acil Çağrı: BLE Köprü vs.\ GSM}
    \begin{columns}
    \column{0.5\textwidth}
        \begin{block}{GSM modülü (SIM800/7600)}
            \begin{itemize}
                \item +Maliyet: 250+ TL
                \item +SIM kart yönetimi
                \item +Operatör onayı
            \end{itemize}
        \end{block}
    \column{0.5\textwidth}
        \begin{block}{ESP32-CAM BLE + Telefon (seçildi)}
            \begin{itemize}
                \item Sürücünün telefonu zaten orada
                \item Ek BLE modülü gerek yok
                \item iPhone: \texttt{tel://}<numara> dialer
            \end{itemize}
        \end{block}
    \end{columns}

    \vspace{0.4cm}
    iOS sandbox tam-otomatik aramaya izin vermez; dialer'a tek tık
    onayı \alert{yanlış pozitif arama koruması} olarak savunuluyor.
\end{frame}

% =========================================================================
% Bölüm 3 — SİSTEM MİMARİSİ
% =========================================================================
\section{Sistem Mimarisi}

\begin{frame}{Uçtan Uca Mimari}
    \begin{center}
    \fbox{\parbox{0.85\linewidth}{\centering\vspace{4em}\textit{[Şekil yer tutucu --- TikZ figürü derlemede atlandı]}\vspace{4em}}}
    \end{center}

    Üç bölge: Araç içi · Bulut tanıma + auth · Sürücü telefonu
\end{frame}

\begin{frame}{Veri Akışı --- 10-Frame Burst}
    \begin{enumerate}
        \item \textbf{TARA} (B1 USER) basıldı $\rightarrow$ STM SCANNING
        \item STM $\rightarrow$ ESP \texttt{CAPTURE} (USART6)
        \item ESP-CAM 5 FPS × 2 s = 10 JPEG yakalar, EC2'ye HTTP POST
        \item Server her frame'de \textbf{kalite filtresi} (blur, yaw, alan, det)
        \item Geçerli embedding'lerin \textbf{centroid}'i $\rightarrow$
              cosine sim $\geq$ 0.40 $\rightarrow$ MATCH
        \item ESP $\rightarrow$ STM \texttt{RESULT:1;name;sim}
        \item STM $\rightarrow$ ESP $\rightarrow$ BLE notify $\rightarrow$ iPhone \texttt{MATCH:..}
        \item iPhone: \texttt{tel://}<numara> dialer ekranı açılır,
              sürücü tek dokunuşla aramayı tamamlar.
    \end{enumerate}

    \vspace{0.3cm}
    Toplam: \alert{$\sim$5--7~s}
\end{frame}

% =========================================================================
% Bölüm 4 — DONANIM ve YAZILIM GERÇEKLEMESİ
% =========================================================================
\section{Gerçekleme}

\begin{frame}{Donanım Listesi}
    \small
    \begin{tabularx}{\linewidth}{X c c}
    \toprule
    \textbf{Parça} & \textbf{Adet} & \textbf{Fiyat}\\
    \midrule
    STM32 NUCLEO-F767ZI                & 1 & (Mevcut)\\
    ESP32-CAM (AI-Thinker, OV3660)     & 1 & 120 TL\\
    ESP32-CAM-MB programlama dock'u    & 1 & 80 TL\\
    Aktif buzzer 5\,V                  & 1 & 15 TL\\
    Push button (PANİK, opsiyonel)     & 1 & 5 TL\\
    Jumper + breadboard                & --- & 60 TL\\
    \midrule
    \textbf{Toplam ek donanım}         &   & \textbf{$\approx$ 280 TL} \\
    \bottomrule
    \end{tabularx}

    \vspace{0.4cm}
    İlk Plan B'deki HM-10 modülü mimariden çıktı; ESP32-CAM aynı işi
    entegre BLE ile yapıyor. Maliyet ve montaj basitleşti.
\end{frame}

\begin{frame}{Sunucu --- Flask + InsightFace + Auth}
    \begin{itemize}
        \item Python 3.11 · Flask · Gunicorn (-w 1) · Docker
        \item AWS EC2 m7i-flex.large · Frankfurt
        \item Tanıma: \texttt{POST /search}
              (10-frame burst, X-Session-Id, JPEG)
        \item \textbf{Kullanıcı yönetimi}: \texttt{/auth/\{register,login,me,logout\}}
            \begin{itemize}
                \item SQLite (\texttt{/app/data/users.db})
                \item werkzeug.security şifre özeti (PBKDF2 + tuz)
                \item \texttt{secrets.token\_urlsafe(32)} bearer token
            \end{itemize}
    \end{itemize}

    \vspace{0.3cm}
    \begin{exampleblock}{Gunicorn -w 1 zorunluluğu}
        Oturum state'i in-memory per-process tutuluyor. 2+ worker burst'ün
        karelerini farklı süreçlere dağıtır $\rightarrow$ agregasyon kırılır.
    \end{exampleblock}
\end{frame}

\begin{frame}{STM32 Olay Yöneticisi}
    \begin{columns}
    \column{0.55\textwidth}
        \begin{itemize}
            \item HAL'dan soyutlanmış \textbf{taşınabilir C99}
            \item Callback vtable: LED, buzzer, UART, clock
            \item 6 durum: IDLE / SCANNING / MATCH / NOMATCH / PANIC / NETERR
            \item PANIC her zaman armed
            \item \alert{HSI $\times$ PLL $\rightarrow$ 216 MHz}
                  (HSE\_VALUE makro hatası bypass'landı)
            \item USART6 (Arduino D0/D1) ESP-CAM köprüsü
        \end{itemize}
    \column{0.45\textwidth}
        \begin{exampleblock}{Host-side birim test}
            \texttt{gcc -O2 -Wall} ile derlenir; \\
            12 senaryo deterministik geçer.\\[2pt]
            Donanım üzerinde flash--flash döngüsü minimize edildi.
        \end{exampleblock}
    \end{columns}
\end{frame}

\begin{frame}{ESP32-CAM --- Kamera + Wi-Fi + BLE Çevresel}
    \begin{itemize}
        \item AI-Thinker, OV3660 sensörü · ESP-WROOM-32 SoC
        \item PlatformIO + Arduino-ESP32 firmware
        \item Wi-Fi STA (sürücü hotspot'u) · HTTP POST → EC2 /search
        \item BLE çevresel: \texttt{TaxiGuard}, FFE0/FFE1, MTU 247
            \begin{itemize}
                \item Varsayılan MTU 23, \texttt{MATCH:...} satırını 20B'da kırıyordu
            \end{itemize}
        \item STM yön mesajlarını (\texttt{SCANNING/MATCH:.../NOMATCH/PANIC/NETERR})
              USART'tan alır, doğrudan BLE notify'a aktarır.
    \end{itemize}

    \vspace{0.3cm}
    \begin{alertblock}{Güç notu}
        Kamera + Wi-Fi + BLE eşzamanlı, $\sim$500 mA pik. STM 5V pininden
        besleme \emph{brown-out}'a yol açar; ayrı USB kaynak şart.
    \end{alertblock}
\end{frame}

% =========================================================================
% Bölüm 5 — iPHONE UYGULAMASI + AUTH
% =========================================================================
\section{Sürücü Uygulaması}

\begin{frame}{iPhone --- TaksiGuvenlik}
    Swift 5.9 + SwiftUI + \texttt{@Observable}, iOS 17+
    \begin{itemize}
        \item \textbf{AuthManager}: \texttt{/auth} REST + Keychain token
        \item \textbf{BLEManager}: CoreBluetooth central, \texttt{TaxiGuard}'a otomatik bağlan
        \item \textbf{RootView}: auth durumuna göre Login / Register / Home (TabView)
        \item \textbf{TabView}:
            \begin{itemize}
                \item Ana Sayfa --- selamlama + plaka + sistem durumu
                \item Tarama --- canlı state + olay günlüğü
                \item Profil --- kullanıcı bilgisi + Çıkış Yap
            \end{itemize}
        \item MATCH/PANIC $\rightarrow$ \texttt{UIApplication.shared.open("tel://\dots")}
    \end{itemize}
\end{frame}

\begin{frame}{Kullanıcı Yönetimi Akışı}
    \begin{center}
    \fbox{\parbox{0.85\linewidth}{\centering\vspace{4em}\textit{[Şekil yer tutucu --- TikZ figürü derlemede atlandı]}\vspace{4em}}}
    \end{center}

    \vspace{0.2cm}
    Plaka regex \texttt{34 ABC 1234} formatında doğrulanır; kayıt başarılı
    olduğunda otomatik login + ana ekran.
\end{frame}

% =========================================================================
% Bölüm 6 — ÖLÇÜM ve SONUÇLAR
% =========================================================================
\section{Ölçüm ve Sonuçlar}

\begin{frame}{Ölçüm Harness'ı}
    \texttt{eval/far\_frr.py}:
    \begin{itemize}
        \item Bir veri seti klasörü (\texttt{<isim>/<foto>.jpg}) tarar
        \item Her fotoğraf 10-frame burst olarak \texttt{/search}'e gönderilir
        \item Eşik $\tau \in [0, 1]$ üzerinde sweep
        \item EER + en iyi operating point (FAR $\leq$ 1\%) raporlanır
        \item ROC eğrisi PNG olarak kaydedilir
    \end{itemize}

    \vspace{0.3cm}
    Çıktılar tezin Tablo 5.x'larına doğrudan kaynak olur.
\end{frame}

\begin{frame}{Ölçüm Planı --- Üç Eksen}
    \begin{enumerate}
        \item \textbf{Aydınlatma}:
              100~lx (alacakaranlık) · 300~lx (ofis) · 600~lx (gün ışığı)
        \item \textbf{Mesafe}:
              30--60~cm · 60--90~cm · 90--120~cm
        \item \textbf{Pasif canlılık}:
              gerçek yüz vs.\ ekran replay
    \end{enumerate}

    \vspace{0.3cm}
    Test seti: ekip + gönüllüler (5--8 kişi).
    Sergi takvimi içinde uygulanabilir kapsam.
\end{frame}

\begin{frame}{Mevcut Durum}
    \small
    \begin{tabular}{l l}
    \toprule
    \textbf{Aşama} & \textbf{Durum}\\
    \midrule
    EC2 sunucu (\texttt{/search} + \texttt{/auth})  & \alert{Canlı, smoke OK}\\
    ESP32-CAM Wi-Fi + HTTP POST + JSON              & \alert{Test edildi}\\
    ESP32-CAM BLE \texttt{TaxiGuard} advertising    & \alert{Test edildi}\\
    STM32 HSI$\times$PLL 216 MHz + USART6           & \alert{Test edildi}\\
    STM32 state machine (host-side)                 & \alert{12/12 test geçti}\\
    STM $\leftrightarrow$ ESP UART köprüsü          & \alert{Test edildi}\\
    iPhone Login/Register + TabView + Keychain      & \alert{Yazıldı, build bekliyor}\\
    iPhone otomatik dialer (\texttt{tel://})         & \alert{Test edildi}\\
    \hline
    FAR/FRR ölçümü                                  & Sergi haftası\\
    \texttt{emergencyNumber} 155'e geri çekme       & Demo öncesi\\
    \bottomrule
    \end{tabular}
\end{frame}

% =========================================================================
% Bölüm 7 — SONUÇ
% =========================================================================
\section{Sonuç}

\begin{frame}{Özgün Katkılar}
    \begin{itemize}
        \item Açık kaynak InsightFace + ESP32-CAM ile düşük maliyetli
              taksi senaryosu entegrasyonu
        \item ESP32-CAM'in entegre BLE çevresel arayüzü ile tek board
              üzerinde kamera + Wi-Fi + BLE; ek modül (HM-10) gereksiz
        \item Cross-frame $\sigma$ ile ek model olmadan pasif canlılık
        \item HAL'dan soyutlanmış, host-side test edilebilen taşınabilir
              STM32 durum makinesi
        \item SQLite + werkzeug + secrets ile ek bağımlılık gerektirmeyen
              hafif kullanıcı kaydı altyapısı
        \item Otomatik FAR/FRR ölçüm aracı $\rightarrow$ tekrarlanabilir tablolar
    \end{itemize}
\end{frame}

\begin{frame}{Gelecek Çalışmalar}
    \begin{enumerate}
        \item \textbf{Faz 2}: Caddy + Let's Encrypt $\rightarrow$ HTTPS;
              \texttt{SHARED\_SECRET} aktif; token süresi (~1/2 gün iş)
        \item \textbf{Android istemcisi}: \texttt{Intent.ACTION\_CALL}
              ile tek-dokunuşsuz çağrı + 5\,s ``İptal'' geri sayımı
        \item \textbf{Edge inference}: ESP32-S3 üzerinde nicemlenmiş ön-eleme
        \item \textbf{Filo entegrasyonu}: MQTT/WebSocket olay panosu, plaka
              eşleme paneli
        \item \textbf{Aktif canlılık}: arka koltukta sessiz kafa-jest doğrulaması
        \item \textbf{Yasal çerçeve}: KVKK açık rıza prosedürü + yetkili
              veri tabanı erişimi
    \end{enumerate}
\end{frame}

\begin{frame}[plain]
    \begin{center}
        \vspace{1.2cm}
        {\Huge \bfseries Teşekkürler}\\[1cm]
        {\Large Sorular?}\\[2cm]
        \small Ekin Ağaoğlu \\
        ekin.agaoglu@example.com \\[6pt]  % TODO: gerçek e-posta
        Danışman: Prof.\ Dr.\ Aydoğan Savran
    \end{center}
\end{frame}

\end{document}
